\documentclass[11pt,letterpaper]{article}
\usepackage[margin=1in]{geometry}
\usepackage{amsmath,amssymb,bm}
\usepackage{booktabs}
\usepackage{enumitem}
\usepackage{xcolor}
\usepackage[colorlinks=true,linkcolor=blue!70!black,urlcolor=blue!70!black,citecolor=blue!70!black]{hyperref}
\usepackage{fancyhdr}
\usepackage{titlesec}
\usepackage{mathtools}

\pagestyle{fancy}
\fancyhf{}
\rhead{Simulating IMU Data for the SO(3) Filter}
\lhead{}
\rfoot{Page \thepage}
\setlength{\headheight}{14pt}

\setlength{\parskip}{5pt}
\setlength{\parindent}{0pt}

\titlespacing*{\section}{0pt}{14pt}{6pt}
\titlespacing*{\subsection}{0pt}{10pt}{4pt}
\titlespacing*{\subsubsection}{0pt}{8pt}{3pt}

% Convenience commands
\newcommand{\R}{\mathbb{R}}
\newcommand{\SO}{\mathrm{SO}(3)}
\newcommand{\expmap}{\mathrm{Exp}}
\newcommand{\eye}{\mathbf{I}}

\title{\textbf{Simulating IMU Data for the SO(3) Attitude Filter}\\[8pt]
\large What We Need, Why We Need It, and How to Generate It}
\author{}
\date{February 2026}

\begin{document}
\maketitle
\thispagestyle{fancy}

\tableofcontents
\newpage

%======================================================================
\section{The Goal}
%======================================================================

We want to test the SO(3) error-state EKF developed in the companion
document.  To test a filter, we need:

\begin{enumerate}[nosep]
  \item \textbf{Ground truth} --- the true orientation $R(t)$ at
    every time step, so we can measure filter error.
  \item \textbf{Synthetic measurements} --- gyroscope and
    accelerometer readings generated from the ground truth, with
    realistic noise and bias, so the filter has something to
    process.
\end{enumerate}

This document explains exactly what ``ground truth'' means for
attitude estimation, what physical assumptions underlie the sensor
models, and how to generate everything in MATLAB.

%======================================================================
\section{Do We Need to Simulate a Full Trajectory?}
%======================================================================

\textbf{Short answer: No.}  We need a \emph{rotation trajectory},
not a translational trajectory.

\subsection{What the Filter Estimates}

The SO(3) filter estimates \emph{orientation only} --- the rotation
matrix $R(t)$ that maps body-frame vectors to world-frame vectors.
It does not estimate position or velocity.  Therefore the ground
truth we need is:
\[
  R_{\mathrm{true}}(t) \in \SO \quad \text{for } t = 0, \Delta t, 2\Delta t, \ldots, T.
\]

We do \emph{not} need to simulate:
\begin{itemize}[nosep]
  \item Position or velocity of the vehicle.
  \item Aerodynamic forces, thrust, or orbital mechanics.
  \item Translational accelerations (with one important caveat ---
    see Section~\ref{sec:accel_assumption}).
\end{itemize}

\subsection{What Defines a Rotation Trajectory}

A rotation trajectory is fully determined by specifying either:
\begin{enumerate}[nosep]
  \item The \textbf{angular velocity} $\bm{\omega}(t) \in \R^{3}$
    (in the body frame) as a function of time, plus an initial
    orientation $R(0)$.
  \item The \textbf{orientation} $R(t)$ directly as a function of
    time (from which $\bm{\omega}(t)$ can be recovered).
\end{enumerate}

Option~1 is simpler for simulation: we choose $\bm{\omega}(t)$,
then propagate $R(t)$ by numerical integration.

%======================================================================
\section{The Accelerometer Assumption}
\label{sec:accel_assumption}
%======================================================================

This is the one place where translational motion matters, and it
deserves careful discussion.

\subsection{What an Accelerometer Actually Measures}

An accelerometer does \emph{not} measure acceleration.  It measures
\textbf{specific force} --- the non-gravitational force per unit
mass acting on the sensor:
\[
  \mathbf{a}_m = R^{\mathsf{T}} (\mathbf{a}_{\mathrm{vehicle}} - \mathbf{g})
    + \mathbf{n}_{a},
\]
where $\mathbf{a}_{\mathrm{vehicle}}$ is the true translational
acceleration of the vehicle in the world frame, and $\mathbf{g}$
is the gravitational acceleration vector.

\subsection{The Static/Quasi-Static Assumption}

In our filter, we use the simplified measurement model:
\[
  \mathbf{a}_m = R^{\mathsf{T}} \mathbf{g} + \mathbf{n}_{a}.
\]

This is valid when $\mathbf{a}_{\mathrm{vehicle}} \approx \mathbf{0}$
--- i.e., the vehicle is either stationary, moving at constant
velocity, or experiencing accelerations that are small compared to
$g = 9.8$\,m/s$^{2}$.

\subsection{When This Assumption Holds}

\begin{center}
\begin{tabular}{lcc}
\toprule
\textbf{Scenario} & \textbf{Typical accel.} &
  \textbf{Assumption valid?} \\
\midrule
Stationary IMU on a table & $0$ & Yes \\
Ground vehicle, normal driving & $< 0.3g$ & Approximately \\
Quadrotor in hover & $\approx 0$ & Yes \\
Quadrotor in aggressive manoeuvre & $1$--$3g$ & No \\
Fixed-wing aircraft, straight \& level & $< 0.1g$ & Yes \\
Fixed-wing aircraft, banking turn & $0.5$--$2g$ & No \\
\bottomrule
\end{tabular}
\end{center}

\subsection{When This Assumption Fails Completely: Satellites}

A satellite in orbit is in \textbf{free fall}.  The gravitational
force and the ``centrifugal force'' (in the orbital frame) cancel
exactly.  An accelerometer on a satellite reads \emph{zero} (plus
noise) --- it provides no attitude information whatsoever.

Satellite attitude determination uses different sensors:
\begin{itemize}[nosep]
  \item \textbf{Star trackers} --- measure the direction to known
    stars in the body frame.
  \item \textbf{Sun sensors} --- measure the direction to the Sun.
  \item \textbf{Magnetometers} --- measure Earth's magnetic field
    vector (LEO only).
  \item \textbf{Horizon sensors} --- measure the direction to
    Earth's limb.
\end{itemize}

The SO(3) filter structure is identical for these sensors --- only
the measurement model $h(R)$ changes.  But for our simulation, we
use the accelerometer model, which means we are simulating a
\textbf{terrestrial or near-surface} scenario.

\subsection{Bottom Line}

For our simulation, we assume the body is not translating (or
translating at constant velocity).  The accelerometer measures the
gravity vector rotated into the body frame, plus noise.  No
translational trajectory is needed.

%======================================================================
\section{Generating the Ground Truth}
%======================================================================

The ground truth consists of three time series: angular velocity,
orientation, and gyro bias.

\subsection{Step 1: Define the True Angular Velocity}

We choose $\bm{\omega}_{\mathrm{true}}(t) \in \R^{3}$ at each
time step.  This is the true body-frame angular velocity.  Several
useful scenarios:

\begin{center}
\begin{tabular}{lp{7cm}}
\toprule
\textbf{Scenario} & \textbf{Description} \\
\midrule
Static &
  $\bm{\omega}(t) = \mathbf{0}$.  Body is fixed at a tilted
  orientation.  Tests bias estimation and steady-state convergence.
  \\[4pt]
Constant rate &
  $\bm{\omega}(t) = (0, 0, \Omega)^{\mathsf{T}}$ (spin about
  $z$-axis).  Tests tracking of constant rotation.
  \\[4pt]
Sinusoidal &
  $\omega_i(t) = A_i \sin(2\pi f_i t + \phi_i)$ for each axis.
  Tests dynamic tracking with smooth, predictable motion.
  \\[4pt]
Multi-axis tumble &
  Superposition of sinusoids at different frequencies and amplitudes
  on all three axes.  Stress test for the filter.
  \\[4pt]
Step input &
  $\bm{\omega}(t)$ steps from zero to a constant value at $t_0$.
  Tests transient response.
  \\
\bottomrule
\end{tabular}
\end{center}

\subsection{Step 2: Propagate the True Orientation}

Given $\bm{\omega}_{\mathrm{true}}(t_k)$ and an initial orientation
$R_{\mathrm{true}}(0)$, propagate forward:
\[
  \boxed{
  R_{\mathrm{true}}(t_{k+1})
  = R_{\mathrm{true}}(t_k) \; \expmap\bigl(\bm{\omega}_{\mathrm{true}}(t_k)
    \, \Delta t\bigr).
  }
\]

This is exact for constant $\bm{\omega}$ over $[t_k, t_{k+1}]$
(the zero-order hold assumption).  For smoothly varying
$\bm{\omega}$, it is accurate to first order in $\Delta t$.  At
100\,Hz, this is more than adequate.

\textbf{Initial orientation:}  Set $R_{\mathrm{true}}(0)$ to
represent the initial tilt.  For example, to start with the body
tilted 30$^{\circ}$ in roll and 20$^{\circ}$ in pitch:
\[
  R_{\mathrm{true}}(0) = R_x(30^{\circ})\, R_y(20^{\circ}).
\]

\subsection{Step 3: Generate the True Gyro Bias}

The gyro bias evolves as a random walk:
\[
  \boxed{
  \mathbf{b}(t_{k+1}) = \mathbf{b}(t_k)
    + \sigma_b \sqrt{\Delta t} \;\mathbf{w}_k,
  \qquad \mathbf{w}_k \sim \mathcal{N}(\mathbf{0}, \eye_3).
  }
\]

\textbf{Initial bias:}  Set $\mathbf{b}(0) = \mathbf{b}_0$ to a
chosen value (e.g., $(0.01, -0.005, 0.008)^{\mathsf{T}}$ rad/s).
This is the ``true'' initial bias that the filter must estimate.

\textbf{Typical values:}  Consumer MEMS gyros have bias instability
of $1$--$10$\,deg/hr, which corresponds to
$\sigma_b \approx 10^{-4}$--$10^{-3}$\,rad/s/$\sqrt{\text{s}}$.
Over a 60-second simulation at $\sigma_b = 5 \times 10^{-4}$, the
bias will random-walk by roughly
$\sigma_b \sqrt{T} = 5 \times 10^{-4} \times \sqrt{60}
\approx 0.004$\,rad/s $\approx 0.2$\,deg/s.  The bias is
``slowly varying'' but definitely not constant.

%======================================================================
\section{Generating the Sensor Measurements}
%======================================================================

With the ground truth in hand, generating measurements is
straightforward.

\subsection{Gyroscope Measurements}

At each time step $t_k$:
\[
  \boxed{
  \bm{\omega}_m(t_k)
  = \bm{\omega}_{\mathrm{true}}(t_k)
    + \mathbf{b}(t_k)
    + \sigma_{\omega} \, \mathbf{v}_k,
  \qquad \mathbf{v}_k \sim \mathcal{N}(\mathbf{0}, \eye_3).
  }
\]

The gyro measurement is the true angular velocity, corrupted by:
\begin{itemize}[nosep]
  \item \textbf{Bias} $\mathbf{b}(t_k)$ --- a slowly drifting
    offset (generated in Step~3 above).
  \item \textbf{White noise}
    $\sigma_{\omega}\,\mathbf{v}_k$ --- zero-mean Gaussian noise
    at each sample, uncorrelated between samples.
\end{itemize}

\textbf{Typical values:}  Consumer MEMS gyros have angular random
walk of $0.1$--$1.0$\,deg/$\sqrt{\text{hr}}$, which converts to
$\sigma_{\omega} \approx 0.003$--$0.03$\,rad/s (at 100\,Hz
sampling).  Tactical-grade gyros are 10--100$\times$ better.

\subsection{Accelerometer Measurements}

At each time step $t_k$ where an accelerometer reading is available:
\[
  \boxed{
  \mathbf{a}_m(t_k)
  = R_{\mathrm{true}}(t_k)^{\mathsf{T}} \, \mathbf{g}
    + \sigma_a \, \mathbf{u}_k,
  \qquad \mathbf{u}_k \sim \mathcal{N}(\mathbf{0}, \eye_3),
  }
\]
where $\mathbf{g} = (0, 0, -9.8)^{\mathsf{T}}$\,m/s$^{2}$.

\textbf{What this means physically:}  The true gravity vector
$(0, 0, -g)$ in the world frame is rotated into the body frame by
$R^{\mathsf{T}}$, then corrupted by white noise.

\textbf{Decimation:}  The accelerometer may run slower than the
gyro.  If the gyro runs at 100\,Hz and the accelerometer at 50\,Hz,
generate accelerometer measurements every other time step.  The
filter runs its prediction step at every gyro sample but only runs
the update step when an accelerometer measurement is available.

\textbf{Typical values:}  Consumer MEMS accelerometers have noise
density of $100$--$500$\,$\mu g/\sqrt{\text{Hz}}$, which at
100\,Hz corresponds to
$\sigma_a \approx 0.01$--$0.05$\,m/s$^{2}$.  For initial testing,
$\sigma_a = 0.5$\,m/s$^{2}$ is a reasonable ``noisy sensor''
value.

\subsection{Accelerometer Bias (Optional Extension)}

In real sensors, accelerometers also have bias.  We omit this for
the initial implementation because:
\begin{itemize}[nosep]
  \item It adds 3 more states to the filter
    ($\delta\mathbf{x} \in \R^{9}$ instead of $\R^{6}$).
  \item Accelerometer bias is harder to observe than gyro bias.
  \item The core concepts are the same --- adding it later is
    straightforward.
\end{itemize}

For now, we model the accelerometer as having white noise only.

%======================================================================
\section{Summary: The Complete Simulation Pipeline}
%======================================================================

\begin{center}
\begin{tabular}{rp{10cm}}
\toprule
\textbf{Step} & \textbf{Action} \\
\midrule
0 & \textbf{Configure:} Choose scenario, sensor parameters
    ($\sigma_{\omega}, \sigma_b, \mathbf{b}_0, \sigma_a$),
    timing ($\Delta t$, $T$), initial orientation $R(0)$. \\[4pt]
1 & \textbf{Generate $\bm{\omega}_{\mathrm{true}}(t)$} from the
    chosen scenario (static, sinusoidal, tumble, etc.). \\[4pt]
2 & \textbf{Propagate $R_{\mathrm{true}}(t)$} from $R(0)$ using
    $R_{k+1} = R_k \, \expmap(\bm{\omega}_k \Delta t)$. \\[4pt]
3 & \textbf{Propagate $\mathbf{b}(t)$} via random walk from
    $\mathbf{b}_0$. \\[4pt]
4 & \textbf{Generate $\bm{\omega}_m(t)$} $= \bm{\omega}_{\mathrm{true}}
    + \mathbf{b} + \text{noise}$. \\[4pt]
5 & \textbf{Generate $\mathbf{a}_m(t)$} $= R_{\mathrm{true}}^{\mathsf{T}}
    \mathbf{g} + \text{noise}$. \\[4pt]
6 & \textbf{Package} all time series into an output struct. \\
\bottomrule
\end{tabular}
\end{center}

\bigskip

The output of this pipeline is a struct containing:

\begin{center}
\begin{tabular}{llp{6cm}}
\toprule
\textbf{Field} & \textbf{Size} & \textbf{Description} \\
\midrule
\texttt{t} & $1 \times N$ & Time vector \\
\texttt{omega\_true} & $3 \times N$ & True angular velocity \\
\texttt{R\_true} & $3 \times 3 \times N$ & True rotation matrices \\
\texttt{b\_true} & $3 \times N$ & True gyro bias \\
\texttt{omega\_m} & $3 \times N$ & Gyro measurements \\
\texttt{a\_m} & $3 \times N$ & Accelerometer measurements \\
\texttt{params} & struct & All sensor/scenario parameters \\
\bottomrule
\end{tabular}
\end{center}

This struct is the input to the SO(3) EKF implementation.  The
ground truth fields (\texttt{R\_true}, \texttt{b\_true},
\texttt{omega\_true}) are used only for computing filter error
after the fact --- the filter never sees them.

%======================================================================
\section{Verification Plots}
%======================================================================

Before running the filter, we should plot the simulated data to
confirm it is reasonable:

\begin{enumerate}[nosep]
  \item \textbf{True angular velocity} $\bm{\omega}(t)$ --- does
    the scenario look as intended?
  \item \textbf{True Euler angles} (extracted from $R_{\mathrm{true}}$
    for visualisation) --- does the orientation evolve smoothly?
  \item \textbf{Gyro measurements vs.\ truth} --- is the bias
    offset visible?  Is the noise level reasonable?
  \item \textbf{Bias evolution} --- is it slowly varying as
    expected?
  \item \textbf{Accelerometer measurements} --- does the measured
    gravity vector rotate correctly with orientation?  Does the
    magnitude stay near $g$?
\end{enumerate}

%======================================================================
\section{What Comes Next}
%======================================================================

With the simulation data in hand, the development sequence is:

\begin{enumerate}
  \item \textbf{\texttt{generate\_imu\_data.m}} --- implements
    everything in this document.  Takes a scenario name and
    parameter struct, returns the data struct.

  \item \textbf{\texttt{run\_so3\_ekf.m}} --- implements the
    filter algorithm from the companion document.  Takes the
    data struct, runs predict/update/inject, returns estimated
    $R(t)$ and $\mathbf{b}(t)$.

  \item \textbf{\texttt{analyse\_results.m}} --- computes
    attitude error $\logmap(R_{\mathrm{true}}\,
    R_{\mathrm{est}}^{\mathsf{T}})$, bias error, plots
    error vs.\ time with $\pm 3\sigma$ bounds from the filter
    covariance.  The $3\sigma$ bounds should contain the errors
    --- if they do, the filter is consistent.

  \item \textbf{Iterate:} Try different scenarios, noise levels,
    initial conditions, and tuning parameters.  Break the filter
    (violate assumptions, use wrong noise parameters) and observe
    how it fails.
\end{enumerate}

\end{document}
